% !TeX root = ../../Book1.tex
% 纲要标题：### D.3 微分 Galois 理论简介
% 纲要位置：计划纲要.md，第 940 行。
% BEGIN APPENDIX OUTLINE SYNC
% #### D.3.1 导子与常数域
% #### D.3.2 线性微分方程与基本解矩阵
% #### D.3.3 Picard--Vessiot 扩张与微分自同构
% ---
% END APPENDIX OUTLINE SYNC

\section{微分 Galois 理论简介}
\label{b1:sec:D03}

经典有限 Galois 理论研究多项式根之间的对称性：固定基域的自同构置换这些根，因而得到有限置换群。线性微分方程的解还带有求导关系；研究它们的对称性时，域自同构也必须与求导交换。对一组基本解，这种自同构表现为常数系数的线性变换，所得群却可能是无限群。这里用一个代数基本解和一个超越基本解比较这两种情形；一般存在定理及 Galois 对应留给后续的线性微分理论。

\subsection{导子与常数域}
\label{b1:subsec:D03:01}

\begin{definition}\label{b1:def:differential-field}
设 \(K\) 为特征零的域，\(\partial:K\rightarrow K\) 为加法映射。若对所有 \(a,b\in K\)，映射 \(\partial\) 都满足 Leibniz 法则\footnote{莱布尼茨（Gottfried Wilhelm Leibniz，1646-1716），德国数学家与哲学家，独立创立微积分并建立其常用记号。}
\[
\partial(ab)=a\partial b+b\partial a,
\]
则称 \((K,\partial)\) 为\term[weifen-yu]{微分域}[differential field]，称 \(\partial\) 为\term[daozi]{导子}[derivation]，称子域 \(C_K=\ker\partial\) 为其\term[changshu-yu]{常数域}[field of constants]。若域扩张上的大域导子限制为小域导子，则称该扩张为微分域扩张。
\end{definition}

这里的导子是域上的一种代数运算，不要求域的元素先是某个区间上的函数。定义中的加法映射只要求 \(\partial(a+b)=\partial a+\partial b\)，没有要求 \(K\)-线性。由 \(\partial(1\cdot1)=2\partial1\) 得 \(\partial1=0\)；若导子还是 \(K\)-线性的，就会有 \(\partial a=a\partial1=0\)，只能得到零导子。对 \(b\neq0\)，在 \(bb^{-1}=1\) 两边施加导子，得到 \(\partial(b^{-1})=-b^{-2}\partial b\)。因此若 \(a,b\) 的导数为零，它们的和、差、积以及 \(b\neq0\) 时的商也如此；这解释了为什么所有导数为零的元素确实组成一个子域。对 \(c\in C_K\)，Leibniz 法则给出 \(\partial(ca)=c\partial a\)，所以导子对常数域 \(C_K\) 线性。所谓常数由导子决定，不能只凭某个元素的记号来判断。

\ProseParagraphBreak
例如设 \(C\) 为特征零的域，在有理函数域 \(C(t)\) 上令 \(\partial t=1\)、\(\partial c=0\)（\(c\in C\)）。加法法则与 Leibniz 法则迫使多项式按通常公式求导，再由商法则确定所有有理函数的导数。它的常数域恰为 \(C\)。事实上，设非零有理函数 \(P/Q\) 已约成互素多项式之比，且导数为零，则 \(P'Q=PQ'\)。于是 \(P\mid P'Q\)，而 \(P,Q\) 在 PID \(C[t]\) 中互素；由 Euclid 引理，\(P\mid P'\)。若 \(P\) 非常数，特征零保证 \(P'\neq0\) 且次数较小，矛盾。所以 \(P\) 为非零常数，原等式进而给出 \(Q'=0\)，故 \(Q\) 也为常数。零函数当然也在 \(C\) 中。

\subsection{线性微分方程与基本解矩阵}
\label{b1:subsec:D03:02}

设 \(K\) 为微分域，\(n\geq1\) 为整数，\(A\in M_n(K)\)。考虑线性方程组 \(\partial z=Az\)，其中未知量 \(z\) 是一个 \(n\) 维列向量，导子逐个作用于分量。若 \(z,w\) 是解而 \(c,d\) 属于常数域，则 \(cz+dw\) 仍是解。系数不能任取于 \(K\)，因为 \(\partial(cz)=(\partial c)z+c\partial z\) 会多出一项。于是解首先组成常数域上的向量空间。

\ProseParagraphBreak
设 \(K\) 为微分域，\(n\geq1\) 为整数，\(A\in M_n(K)\)，\(L/K\) 为微分域扩张。若矩阵 \(Y\in\operatorname{GL}_n(L)\) 满足 \(\partial Y=AY\)，即各列是在 \(L^n\) 中 \(L\)-线性无关的解向量，则称 \(Y\) 为方程 \(\partial z=Az\) 的\term[jiben-jie-juzhen]{基本解矩阵}[fundamental solution matrix]。这里的可逆性是通常的 \(L\)-线性代数条件，解空间的标量域却是 \(C_L\)。对任一解向量 \(z\in L^n\)，可逆性先给出唯一的坐标列 \(c=Y^{-1}z\in L^n\)；还须证明这些坐标全是常数，才能把 \(Y\) 的列看成解空间的一组 \(C_L\)-基。对 \(Y^{-1}Y=I\) 求导，得到
\[
\partial(Y^{-1})Y+Y^{-1}\partial Y=0,
\qquad
\partial(Y^{-1})=-Y^{-1}(\partial Y)Y^{-1}.
\]
于是
\[
\partial(Y^{-1}z)=-Y^{-1}A z+Y^{-1}\partial z=0.
\]
所以 \(c\in C_L^n\)，而 \(z=Yc\)。反过来，每个这样的常数线性组合仍是解。基本解矩阵因而同时记录了全部解以及它们在常数域上的坐标。

\subsection{Picard--Vessiot 扩张与微分自同构}
\label{b1:subsec:D03:03}

微分域的定义并不要求常数域代数闭，例如 \((\mathbb Q(t),\mathrm{d}/\mathrm{d}t)\) 的常数域就是 \(\mathbb Q\)。以下讨论 Picard--Vessiot 扩张时增加常数域代数闭的假设；它是通常存在唯一性及对应定理所采用的条件。

\begin{definition}\label{b1:def:picard-vessiot}
设 \(K\) 为微分域，其常数域 \(C_K=C\) 代数闭，\(n\geq1\) 为整数，并取 \(A\in M_n(K)\)。若微分域扩张 \(L/K\) 满足以下条件：存在可逆矩阵 \(Y\in\operatorname{GL}_n(L)\) 使 \(\partial Y=AY\)；域 \(L\) 由 \(Y\) 的矩阵元在 \(K\) 上生成；且 \(C_L=C\)，则称 \(L/K\) 为方程 \(\partial Y=AY\) 的\term[Picard-Vessiot-kuozhang]{Picard--Vessiot 扩张}[Picard--Vessiot extension]。\footnote{皮卡尔（Charles Émile Picard，1856-1941），法国数学家，研究函数论与微分方程。}\footnote{韦西奥（Ernest Vessiot，1865-1952），法国数学家，发展线性微分方程的 Galois 理论。}集合
\[
\operatorname{Gal}_{\partial}(L/K)=
\{\sigma\in\operatorname{Aut}_K(L)\mid\sigma\partial=\partial\sigma\}
\]
在复合下构成群，称为该扩张的\term[weifen-galois-qun]{微分 Galois 群}[differential Galois group]。
\end{definition}

基本解矩阵存在保证 \(L\) 中已经容纳一组完整的独立解；生成条件则要求只加入这些解及其域运算所需的元素，而不再加入无关的域元素。既然 \(L\) 是域，\((\det Y)^{-1}\) 已经在其中。“不增加常数”排除额外的常数扩张，使解之间的系数始终取自原来的 \(C\)。例如方程 \(y'=0\) 在 \(K\) 中已有基本解 \(1\)，并不需要为它另添常数。这些条件是对给定扩张的检验；定义本身没有证明每个方程都存在这样的扩张。

\ProseParagraphBreak
若 \(\sigma\) 固定 \(K\) 并与导子交换，它便把 \(\partial Y=AY\) 变成 \(\partial\sigma(Y)=A\cdot\sigma(Y)\)。因此它将一组基本解变成另一组基本解。\(\sigma(Y)\) 的每一列都是 \(Y\) 各列的常数线性组合，把这些系数依次放成矩阵的各列，便得到右乘形式 \(\sigma(Y)=YM_\sigma\)。

\begin{proposition}\label{b1:prop:differential-galois-matrices}
设 \(K\) 为微分域，其常数域 \(C=C_K\) 为代数闭域，\(n\geq1\) 为整数；取 \(A\in M_n(K)\)，并设 \(L/K\) 为方程 \(\partial Y=AY\) 的 Picard--Vessiot 扩张，其中 \(Y\in\operatorname{GL}_n(L)\) 是生成 \(L\) 的一个基本解矩阵。则每个 \(\sigma\in\operatorname{Gal}_{\partial}(L/K)\) 唯一写成 \(\sigma(Y)=YM_\sigma\)，其中 \(M_\sigma\in\operatorname{GL}_n(C)\)，并且 \(\sigma\mapsto M_\sigma\) 是单同态。
\end{proposition}

\begin{proof}
因为 \(Y\) 和 \(\sigma(Y)\) 可逆，等式 \(\sigma(Y)=YM_\sigma\) 唯一决定可逆矩阵 \(M_\sigma=Y^{-1}\cdot\sigma(Y)\)。利用前面得到的逆矩阵求导公式，再使用 \(\sigma(A)=A\) 与微分交换条件，得
\[
\partial M_\sigma=-Y^{-1}A\cdot\sigma(Y)+Y^{-1}A\cdot\sigma(Y)=0.
\]
因此矩阵元属于 \(C_L=C\)。仅由两矩阵可逆，只能知道 \(M_\sigma\) 的矩阵元在 \(L\) 中；与导子交换的条件才通过这个零导数计算把它们限制到常数域。这里 \(C\subseteq K\)，所以每个 \(K\) 自同构固定这些系数，进而
\[
\sigma\tau(Y)=\sigma(YM_\tau)=\sigma(Y)M_\tau
=YM_\sigma M_\tau.
\]
由表示矩阵的唯一性，\(M_{\sigma\tau}=M_\sigma M_\tau\)，得到同态性。若 \(M_\sigma=I\)，则 \(\sigma\) 固定 \(Y\) 的全部矩阵元，也固定它们在 \(K\) 上的所有有理表达式。生成条件给出 \(\sigma=\operatorname{id}_L\)，故同态单射。
\end{proof}

这个单同态不一定满射：一个常数矩阵即使把基本解变成另一组基本解，也还必须保持 \(Y\) 的矩阵元在 \(K\) 上满足的全部代数关系，才可能来自域自同构。下面的 \(u^2=t\) 正好展示这一限制。

\ProseParagraphBreak
这里的唯一性以选定的基本解矩阵 \(Y\) 为前提。若 \(\widetilde Y\in\operatorname{GL}_n(L)\) 是同一方程的另一基本解矩阵，则 \(\partial(Y^{-1}\widetilde Y)=0\)，故 \(\widetilde Y=YP\)，其中 \(P\in\operatorname{GL}_n(C)\)。这就是对解空间作常数域上的换基。由于 \(\sigma\) 固定 \(C\)，相对于 \(\widetilde Y\) 的表示矩阵满足
\[
\sigma(\widetilde Y)=\sigma(Y)P
=\widetilde Y\bigl(P^{-1}M_\sigma P\bigr),
\qquad \widetilde M_\sigma=P^{-1}M_\sigma P.
\]
这与普通群表示的换基公式相同。因此微分自同构群由扩张 \(L/K\) 决定，它在 \(\operatorname{GL}_n(C)\) 中的具体矩阵像则随基本解矩阵的选择而共轭。

\ProseParagraphBreak
取特征零的代数闭域 \(C\)，令 \(K=C(t)\)、\(L=C(u)\)，在 \(K\) 上规定 \(\partial c=0\)（\(c\in C\)）及 \(\partial t=1\)，再把 \(K\) 通过 \(t=u^2\) 嵌入 \(L\)。要求导子限制回 \(K\) 的原导子，就不能任意指定 \(\partial u\)：对关系 \(u^2=t\) 求导，必须有 \(2u\partial u=1\)，所以 \(\partial u=1/(2u)=u/(2t)\)。在 \(C(u)\) 上取导子 \((2u)^{-1}\mathrm{d}/\mathrm{d}u\)，确实满足这一条件，限制到 \(K\) 就是原来的导子。对 \(f\in C(u)\)，因为乘子 \((2u)^{-1}\neq0\)，有
\[
\partial f=0\quad\Longleftrightarrow\quad
\frac{\mathrm{d}f}{\mathrm{d}u}=0.
\]
由前面有理函数域的常数域计算，常数域仍为 \(C\)。

\ProseParagraphBreak
非零元素 \(u\) 满足一阶方程 \(y'=y/(2t)\)，可以看成一个 \(1\times1\) 的可逆解矩阵；\(L=K(u)\)，所以三个 Picard--Vessiot 条件逐一成立。\(t\) 不是 \(C(t)\) 中的平方：平方有理函数在 \(t=0\) 处的零或极点次数为偶数，而 \(t\) 的次数为一。因此 \([L:K]=2\)，两个 \(K\)-自同构是 \(u\mapsto u\) 与 \(u\mapsto-u\)。对非恒等者 \(\sigma\)，有
\[
\partial\sigma(u)=-\frac1{2u}
=\sigma\bigl(\partial u\bigr).
\]
两种复合运算 \(\partial\circ\sigma\) 与 \(\sigma\circ\partial\) 在生成元 \(u\) 上相同，在 \(C\) 上均为零，在 \(t\) 上均取值 \(1\)；由加法、乘法及商法则可知二者在整个 \(L\) 上相同。故微分群恰为 \(\{1,-1\}\subseteq C^\times\)。由基本解矩阵的计算，\(L\) 中所有非零解都是 \(cu\)（\(c\in C^\times\)），但若替换 \(u\mapsto cu\) 来自 \(K\)-自同构，就还要有
\[
(cu)^2=\sigma(u^2)=\sigma(t)=t,
\]
即 \(c^2=1\)，所以只有 \(c=\pm1\) 能实现。解空间上的非零常数比例都保持微分方程，域自同构还须保持代数关系；二者的要求在这里确实不同。这个例子中基本解恰好代数，Picard--Vessiot 的定义并未要求一般基本解在基域上代数。

\begin{example}[超越基本解]\label{b1:exm:D-exponential-extension}
在 \(K=\mathbb C(t)\) 上考虑 \(\partial y=y\)。借助\secref{b1:sec:D01}的 Laurent 级数域，在 \(\mathbb C((t))\) 中取形式求导，并令
\[
e^t\coloneqq\sum_{m\geq0}\dfrac{t^m}{m!},\qquad
y=e^t,\qquad L=K(y).
\]
这里 \(e^t\) 只是所写形式幂级数的记号。逐项形式求导给出 \(\partial y=y\)，而常数项为一保证 \(y\neq0\)；以下论证不需要级数的收敛性。

先证明这个基本解超越于 \(K\)。若它代数，取首一最小多项式
\[
F(X)=X^n+\sum_{j=0}^{n-1}a_jX^j\in K[X].
\]
因为 \(y\neq0\)，有 \(a_0\neq0\)：否则 \(F(X)=XH(X)\)，从 \(F(y)=0\) 得 \(H(y)=0\)，与最小次数矛盾。由 \(\partial y=y\)，对 \(F(y)=0\) 求导得到
\[
ny^n+\sum_{j=0}^{n-1}(a_j'+ja_j)y^j=0.
\]
这里 \(a_j'\) 来自系数求导，而 \(ja_jy^j\) 来自 \(\partial(y^j)=jy^j\)。减去 \(nF(y)=0\)，得到
\[
\sum_{j=0}^{n-1}\bigl(a_j'+(j-n)a_j\bigr)y^j=0.
\]
若这些系数不全为零，就得到一个次数小于 \(n\) 的非零 \(K\)-多项式，在 \(y\) 处取值为零，与 \(F\) 的最小次数矛盾。因此各系数全为零，特别是 \(a_0'=na_0\)。写 \(a_0=P/Q\)，其中 \(P,Q\in\mathbb C[t]\) 均非零，便得
\[
P'Q-PQ'=nPQ.
\]
右边因特征零而非零，所以左边也须非零，但其次数至多为 \(\deg P+\deg Q-1\)；右边的次数恰为 \(\deg P+\deg Q\)，矛盾。故 \(y\) 超越。

常数域也可以直接检查。在特征零下，Laurent 级数 \(\sum_{j\geq N}c_jt^j\) 的导数为零，当且仅当所有 \(j\neq0\) 的系数为零。因此 \(\mathbb C((t))\) 的常数域恰为 \(\mathbb C\)；由 \(\mathbb C\subseteq L\subseteq\mathbb C((t))\)，有 \(C_L=\mathbb C\)。于是 \(y\) 是生成 \(L\) 的可逆一阶基本解，且没有增加常数，\(L/K\) 确实是 Picard--Vessiot 扩张。

若 \(\sigma\) 是微分 \(K\) 自同构，则 \(\partial\sigma(y)=\sigma(y)\)，商法则给出
\[
\partial\bigl(\sigma(y)/y\bigr)=0,
\qquad \sigma(y)=cy\quad(c\in\mathbb C^\times).
\]
反过来，\(y\) 超越给出 \(K(y)\cong K(T)\)，其中 \(T\) 是不定元。在有理函数域中，替换 \(T\mapsto cT\) 把 \(P(T)/Q(T)\) 送到 \(P(cT)/Q(cT)\)，非零分母仍非零，且逆替换为 \(T\mapsto c^{-1}T\)。所以每个 \(y\mapsto cy\) 都定义一个 \(K\) 自同构；它在 \(t,y\) 上与导子交换，故由域运算法则在整个 \(L\) 上交换。于是
\[
\operatorname{Gal}_{\partial}(L/K)\cong\mathbb C^\times
=\operatorname{GL}_1(\mathbb C).
\]
前一个例子的代数关系 \(u^2=t\) 把比例限制为 \(\pm1\)，这里的基本解在 \(K\) 上没有非零多项式关系，所有非零常数比例都能实现，所得群便是无限群。这两例说明，确定微分 Galois 群时，既要知道基本解允许怎样作常数线性变换，也要检验这些变换是否保持它们在基域上的代数关系。
\end{example}

一般理论进一步证明，微分 Galois 群的矩阵像是 \(\operatorname{GL}_n(C)\) 的 Zariski 闭子群。\footnote{扎里斯基（Oscar Zariski，1899-1986），俄国出生的美国数学家，以交换代数方法奠定现代代数几何的基础，研究奇点与局部性质。}这表示它是在可逆矩阵中由矩阵元满足的一组 \(C\)-系数多项式方程切出的子群；这些方程反映了保持基本解全部代数关系的要求。上面第一个群在 \(\operatorname{GL}_1(C)=C^\times\) 中由 \(c^2-1=0\) 切出，第二个群则等于整个 \(\operatorname{GL}_1(\mathbb C)\)，没有额外方程限制。相应的微分 Galois 对应把 \(K\) 与 \(L\) 之间的微分子域同群的 Zariski 闭子群对应起来。闭性及对应定理还需专门证明，不能仅凭本节的单同态得出；详细陈述见 Singer 的作者公开稿\cite[Definition 1.3.5, Proposition 1.3.6 and Theorem 1.3.9]{Singer2009DifferentialGalois}。

\ProseParagraphBreak
这里的闭性与第二十二章的 Krull 拓扑闭性有不同的含义。Krull 拓扑以自同构在有限子扩张上的限制为依据：一个自同构属于子群的闭包，当且仅当在每个有限子扩张上都能找到该子群的元素与它有相同限制。Zariski 拓扑则以矩阵元的多项式方程为依据，把所有必须满足的代数关系同时保留下来。第三卷附录 E 引入微分代数对象，第五卷附录 E 讨论线性微分方程及其群。
% ---
